\tinysidebar{\debug{exercises/{test/question.tex}}}
\begin{ex}
Here's something from wikipedia.org:
\begin{console}[commandchars=\~\@\$]
~textbf@1729$ is known as the ~textbf@Hardy-Ramanujan number$, after a
famous anecdote of the British mathematician
~href@http://en.wikipedia.org/wiki/G._H._Hardy$@G. H. Hardy$ regarding a
hospital visit to the Indian mathematician
~href@http://en.wikipedia.org/wiki/Srinivasa_Ramanujan$@Srinivasa
Ramanujan$. In Hardy's words:

~lq~lq I remember once going to see him when he was ill at
~href@http://en.wikipedia.org/wiki/Putney$@Putney$. I had ridden in taxi cab number 1729
and remarked that the number seemed to me rather a dull
one, and that I hoped it was not an unfavorable ~href@http://en.wikipedia.org/wiki/Omen$@omen$.
"No," he replied, "it is a very interesting number;
~textbf@it is the smallest number expressible as the sum of two$
~textbf@cubes in two different ways."$

The quotation is sometimes expressed using the
term "positive cubes", as the admission of negative
perfect cubes (the cube of a ~href@http://en.wikipedia.org/wiki/Negative_and_non-negative_numbers$@negative$ ~href@http://en.wikipedia.org/wiki/Integer$@integer$) gives
the smallest solution as ~href@http://en.wikipedia.org/wiki/91_%28number%29$@91$ (which is a factor of 1729):

~tab[3em]@91 = 6^3 + (-5)^3 =$
~tab[3em]@4^3 + 3^3$

Of course, equating "smallest" with "most negative", as
opposed to "closest to zero" gives rise to solutions
like -91, -189, -1729, and further negative numbers.
This ambiguity is eliminated by the term "positive cubes".

Numbers such as

~tab[3em]@1729 = 1^3 + 12^3 =$
~tab[3em]@9^3 + 10^3$

that are the smallest number that can be expressed as the
sum of two cubes in ~texttt@n$ distinct ways have been dubbed
~href@http://en.wikipedia.org/wiki/Taxicab_number$@taxicab numbers$. 1729 is the second taxicab number (the
first is 2 = 1^3 + 1^3). The number was also found in one
of Ramanujan's notebooks dated years before the incident.
\end{console}
Find all positive integer solutions to the equation
\begin{center}
x$^3$ + y$^3$ = 1729
\end{center}
(Hint: Write a double for-loop.)
\end{ex}
